% 第5章 考察
% 対応する草稿: research/thesis-drafts/chapter5-discussion-skeleton.md

\subsection{結果の解釈}\label{subsec:interpretation}

% TODO: RQ ごとの考察

\subsection{先行研究との対比}\label{subsec:comparison}

% TODO: Kadavath 2022, Tian 2023 との比較

\subsection{DEL 公理との対応}\label{subsec:del-axioms}

% TODO: 真理性公理の充足度（ECE）
% 負の内省（AUROC）
% LLM はどの論理体系に近いか
% DEL-ToM との関係

\subsection{実用的含意}\label{subsec:implications}

% TODO: 日本語LLM利用者への指針

\subsection{研究の限界}\label{subsec:limitations}

% TODO: サンプル数、時間的有効性、タスク多様性
